% ============================================================
% RESULTS — target: 800--1200 words
% Cardinal rule: REPORT ONLY. No interpretation.
% ============================================================

\section{Results}

\subsection{Classifier Validation}

[TODO: Report Cohen's $\kappa$ and per-stage accuracy on 200 hand-labeled examples.
Report Stage~1 (pattern matching) precision/recall and Stage~2 (LLM judge) accuracy
on the subset that passes through.]

\subsection{Baseline Accuracy}

To establish model capability independent of hint injection, we first report accuracy
on the 500 unhinted questions.

% Table 2: Baseline accuracy by model
% \begin{table}[H]
%   \centering
%   \caption{Baseline accuracy (\%) on MMLU (300 questions) and GPQA Diamond
%   (200 questions) without hint injection. Models sorted by overall accuracy.}
%   \label{tab:baseline}
%   \begin{tabular}{lccc}
%     \toprule
%     \textbf{Model} & \textbf{MMLU} & \textbf{GPQA} & \textbf{Overall} \\
%     \midrule
%     [TODO: Fill after experiments] \\
%     \bottomrule
%   \end{tabular}
% \end{table}

[TODO: Report baseline accuracy results. Reference Table~\ref{tab:baseline}.]

\subsection{Hint Influence Rates}

Before measuring faithfulness, we report the rate at which each hint type
successfully influenced model answers.

% Figure 1: Hint influence rate heatmap (model x hint type)
% \begin{figure}[H]
%   \centering
%   \includegraphics[width=\textwidth]{fig1_influence_heatmap.pdf}
%   \caption{Hint influence rate (\%) by model and hint type. Darker cells
%   indicate higher susceptibility to the hint. [KEY TAKEAWAY].}
%   \label{fig:influence}
% \end{figure}

[TODO: Report influence rates. Which models are most/least susceptible?
Which hint types are most effective? Reference Figure~\ref{fig:influence}.]

\subsection{Faithfulness Rates by Model and Hint Type (H1)}

To examine whether faithfulness rates vary across model families (H1), we
computed faithfulness rates for each model-hint combination.

% Figure 2: Main result — faithfulness rate heatmap (model x hint type)
% \begin{figure}[H]
%   \centering
%   \includegraphics[width=\textwidth]{fig2_faithfulness_heatmap.pdf}
%   \caption{Faithfulness rate (\%) by model and hint type. Values represent
%   the proportion of hint-influenced responses where the CoT explicitly
%   acknowledges the hint. 95\% bootstrap CIs in parentheses.
%   [KEY TAKEAWAY].}
%   \label{fig:faithfulness}
% \end{figure}

% Table 3: Full faithfulness results with CIs
% \begin{table}[H]
%   \centering
%   \caption{Faithfulness rates (\%) with 95\% bootstrap confidence intervals.
%   Bold indicates highest faithfulness per hint type.}
%   \label{tab:faithfulness}
%   [TODO: 12 rows x 7 columns (model + 6 hint types + overall)]
% \end{table}

[TODO: Report faithfulness rates. Reference Figure~\ref{fig:faithfulness} and
Table~\ref{tab:faithfulness}. Report McNemar's test results for pairwise
model comparisons. State whether H1 is supported.]

\subsection{Hint-Type Patterns (H2)}

To examine whether certain hint types exhibit consistently lower faithfulness
across all models (H2), we compared faithfulness rates across the six hint
categories.

% Figure 3: Faithfulness by hint type (box plots across models)
% \begin{figure}[H]
%   \centering
%   \includegraphics[width=0.8\textwidth]{fig3_faithfulness_by_hint.pdf}
%   \caption{Distribution of faithfulness rates across 12 models for each
%   hint type. Box plots show median, IQR, and range. [KEY TAKEAWAY].}
%   \label{fig:by_hint}
% \end{figure}

[TODO: Report per-hint-type results. Which hints are most faithfully
reported? Which are systematically underreported? State whether H2 is
supported.]

\subsection{Scaling Analysis}

To examine the relationship between model scale and faithfulness, we plot
faithfulness rates against active parameter count.

% Figure 4: Faithfulness vs. scale (scatter with regression line)
% \begin{figure}[H]
%   \centering
%   \includegraphics[width=0.8\textwidth]{fig4_faithfulness_vs_scale.pdf}
%   \caption{Faithfulness rate vs.\ active parameter count (log scale).
%   Each point is one model; error bars show 95\% CIs. Dashed line shows
%   log-linear regression ($R^2 = $ [TODO], $p = $ [TODO]).}
%   \label{fig:scaling}
% \end{figure}

[TODO: Report scaling analysis. Is there a correlation between model size and
faithfulness? Reference Figure~\ref{fig:scaling}. Report regression
statistics.]

\subsection{CoT Length Analysis}

Following Chen et al.~\cite{chen2025reasoning}, we compare the length of faithful
vs.\ unfaithful CoTs.

% Figure 5: CoT length distribution (faithful vs. unfaithful)
% \begin{figure}[H]
%   \centering
%   \includegraphics[width=0.8\textwidth]{fig5_cot_length.pdf}
%   \caption{Distribution of CoT token counts for faithful (hint acknowledged)
%   and unfaithful (hint not acknowledged) responses, pooled across all
%   models. [KEY TAKEAWAY regarding length differences].}
%   \label{fig:cot_length}
% \end{figure}

[TODO: Report whether unfaithful CoTs are longer, replicating the finding
of Chen et al. Reference Figure~\ref{fig:cot_length}. Report median
lengths and statistical test.]

\subsection{Replication of Chen et al.}

To validate our methodology, we compare our DeepSeek-R1 results with those
reported by Chen et al.~\cite{chen2025reasoning}.

[TODO: Report side-by-side comparison of our DeepSeek-R1 faithfulness rates
vs.\ Chen et al.'s reported 39\%. Discuss any discrepancies and their
likely causes.]

\subsection{Summary}

[TODO: ``In summary, results [supported/did not support] H1 and [supported/did
not support] H2. The key findings are: (1)~[FINDING], (2)~[FINDING],
(3)~[FINDING]. These results are examined in the following section.'']
